\scanpage{58}
\mspage{14}

C) \uline{Découpage} \deleted{\unsure{de|des}} suivant une partie \inserted{ouv. et f.} $\Sigma'\subset\Sigma_0$ de la sous-variété fermée des cercles de découpage -- cela correspond à une partie $A'$ de $A_0$ (arêtes \deleted{\unsure{ord.}}\inserted{liées}), mais cette fois ci il ne s'agit pas de ``gommer'' i.e de contracter en un point, mais de ``découper'', ce qui \unsure{au niveau des graphes|sur nos amis les graphes} revient à découper une arête par son milieu (i.e. à la remplacer par deux arêtes libres).

Le \uline{marquage suit}. \deleted{Ici encore, pour \unsure{un|une} \unsure{a}} Mais il y a ici à distinguer entre le cas d'une opération anodine \inserted{en $C$} -- celle où \deleted{il y a bel et \unsure{bien}} \inserted{$C$ est effectivement marqué} i.e. $C\cap R\ne\varnothing$ i.e. $C\subset\overset{\unsure{\mathrm{cc}}}{\Sigma}^+$ -- \deleted{\unsure{et|est}} à peu de choses près, (une ambiguïté dans un ens. fini, \deleted{\unsure{)|ou}} ou pas) peut se \unsure{réverser|réinverser}, à l'aide du produit de cette opération -- et celui où $C$ n'est pas marqué, Mais il semble que dans la façon dont nous allons travailler

% archive sheet number: 57

